\section{Wprowadzenie i mapa systemów eClinical}

% SLIDE ID: sec01-goal
\BlockGoalFrame{\SlideID{sec01-goal}}{
  \item zobaczyć mapę najważniejszych systemów eClinical,
  \item zrozumieć, jak wspierają przebieg badania klinicznego,
  \item połączyć teorię z demo i pracą warsztatową.
}

% SLIDE ID: sec01-audience
\begin{frame}{Dla kogo jest ten wykład}
  \SlideID{sec01-audience}
  \begin{itemize}
    \item dla osób ze studiów podyplomowych związanych z badaniami klinicznymi,
    \item dla przyszłych koordynatorów, monitorów i członków zespołów projektowych,
    \item dla uczestników, którzy chcą rozumieć systemy używane w praktyce.
  \end{itemize}
\end{frame}

% SLIDE ID: sec01-assumptions
\begin{frame}{Założenia na start}
  \SlideID{sec01-assumptions}
  \begin{itemize}
    \item znacie podstawy badań klinicznych,
    \item nie musicie znać systemów informatycznych,
    \item zaczynamy od logiki procesu, a dopiero potem przechodzimy do narzędzi.
  \end{itemize}
\end{frame}

% SLIDE ID: sec01-agenda
\begin{frame}{Agenda dnia}
  \SlideID{sec01-agenda}
  \small
  \begin{tabularx}{\linewidth}{>{\raggedright\arraybackslash}p{0.10\linewidth} >{\raggedright\arraybackslash}X}
    \toprule
    Blok & Temat \\
    \midrule
    1 & Wprowadzenie i mapa systemów eClinical \\
    2 & EDC / eCRF \\
    3 & Warsztat EDC \\
    4 & Dzienniczek pacjenta / ePRO \\
    5 & IWRS / IVRS \\
    6 & CTMS \\
    7 & eTMF \\
    8 & Safety / AE / SAE \\
    9 & AI w badaniach klinicznych \\
    10 & Case study: jedno badanie od A do Z \\
    11 & Podsumowanie i pytania \\
    \bottomrule
  \end{tabularx}
\end{frame}

% SLIDE ID: sec01-map
\begin{frame}{Mapa systemów eClinical}
  \SlideID{sec01-map}
  \begin{columns}[T,onlytextwidth]
    \column{0.48\textwidth}
    \begin{block}{Systemy widoczne w pracy z danymi}
      \begin{itemize}
        \item EDC / eCRF,
        \item ePRO,
        \item safety.
      \end{itemize}
    \end{block}

    \column{0.48\textwidth}
    \begin{block}{Systemy operacyjne i nadzorcze}
      \begin{itemize}
        \item IWRS / IVRS,
        \item CTMS,
        \item eTMF.
      \end{itemize}
    \end{block}
  \end{columns}
\end{frame}

% SLIDE ID: sec01-story
\begin{frame}{Główna historia dnia}
  \SlideID{sec01-story}
  \begin{itemize}
    \item przejdziemy przez jedno badanie kliniczne od startu do zamknięcia,
    \item w każdym kroku zobaczymy, jaki system wspiera dany proces,
    \item dzięki temu łatwiej będzie połączyć narzędzie z realną pracą zespołu.
  \end{itemize}
\end{frame}

% SLIDE ID: sec01-workflow
\begin{frame}{Jak będziemy pracować}
  \SlideID{sec01-workflow}
  \begin{columns}[T,onlytextwidth]
    \column{0.24\textwidth}
    \begin{beamercolorbox}[wd=\linewidth,sep=0.7em,rounded=true]{block body}
      \centering
      \textbf{Teoria}
    \end{beamercolorbox}

    \column{0.24\textwidth}
    \begin{beamercolorbox}[wd=\linewidth,sep=0.7em,rounded=true]{block body}
      \centering
      \textbf{Demo}
    \end{beamercolorbox}

    \column{0.24\textwidth}
    \begin{beamercolorbox}[wd=\linewidth,sep=0.7em,rounded=true]{block body}
      \centering
      \textbf{Ćwiczenie}
    \end{beamercolorbox}

    \column{0.24\textwidth}
    \begin{beamercolorbox}[wd=\linewidth,sep=0.7em,rounded=true]{block body}
      \centering
      \textbf{Dyskusja}
    \end{beamercolorbox}
  \end{columns}

  \vspace{0.6cm}

  \begin{itemize}
    \item EDC będzie dostępne także dla uczestników,
    \item pozostałe systemy pokazuje prowadzący,
    \item dodatkowo zobaczymy dzienniczek pacjenta / ePRO.
  \end{itemize}
\end{frame}

% SLIDE ID: sec01-discussion
\DiscussionFrame{Z czym uczestnicy wchodzą na ten wykład?}{
  Który z systemów jest dziś dla Was najbardziej nieznany: \texttt{EDC}, \texttt{CTMS}, \texttt{IWRS}, \texttt{eTMF}, safety czy \texttt{ePRO}? W którym miejscu procesu badania najtrudniej było dotąd zrozumieć rolę narzędzi informatycznych?
}
